\textbf{Phần 1. Trắc nghiệm lựa chọn}
\setcounterpageref{ex}{0}

\begin{ex} %[5H1N8]
Số thích hợp điền vào dấu hỏi chấm để $6\text{ yến} = ?\text{ kg}$ là
\choice
{$6$}
{$600$}
{$6\,000$}
{\True $60$}
\loigiai{
Vì $1\text{ yến} = 10\text{ kg}$ nên $6\text{ yến} = 60\text{ kg}$.\\
\textbf{Đáp án của câu hỏi là D}
}
\end{ex}

\begin{ex} %[5H1N8]
Số thích hợp điền vào dấu hỏi chấm để $2\text{ tấn} = ?\text{ kg}$ là
\choice
{\True $2\,000$}
{$20$}
{$200$}
{$20\,000$}
\loigiai{
Vì $1\text{ tấn} = 1\,000\text{ kg}$ nên $2\text{ tấn} = 2\,000\text{ kg}$.\\
\textbf{Đáp án của câu hỏi là A}
}
\end{ex}

\begin{ex} %[5H1H8]
Số thích hợp điền vào chỗ chấm để $1\text{ tấn} = \dots\text{ yến}$ là
\choice
{$10$}
{$1\,000$}
{$10\,000$}
{\True $100$}
\loigiai{
Vì $1\text{ tấn} = 1\,000\text{ kg}$ và $1\text{ yến} = 10\text{ kg}$ nên $1\text{ tấn} = 100\text{ yến}$.\\
\textbf{Đáp án của câu hỏi là D}
}
\end{ex}

\begin{ex} %[5H1H8]
Góc có số đo bằng $120^\circ$ là
\choice
{Góc nhọn}
{\True Góc tù}
{Góc vuông}
{Góc bẹt}
\loigiai{
Góc lớn hơn $90^\circ$ và nhỏ hơn $180^\circ$ là góc tù. Do đó góc $120^\circ$ là góc tù.\\
\textbf{Đáp án của câu hỏi là B}
}
\end{ex}

\begin{ex} %[5H1H8]
Số thích hợp điền vào dấu hỏi chấm để $5\text{ giờ} = ?\text{ phút}$ là
\choice
{$50$}
{$250$}
{\True $300$}
{$360$}
\loigiai{
Vì $1\text{ giờ} = 60\text{ phút}$ nên $5\text{ giờ} = 5 \times 60 = 300\text{ phút}$.\\
\textbf{Đáp án của câu hỏi là C}
}
\end{ex}

\begin{ex} %[5H1H8]
Đổi $\dfrac{1}{10}\text{ giờ}$ ra phút ta được kết quả là
\choice
{$10\text{ phút}$}
{$12\text{ phút}$}
{$15\text{ phút}$}
{\True $6\text{ phút}$}
\loigiai{
Ta có: $\dfrac{1}{10}\text{ giờ} = 60 \times \dfrac{1}{10} = 6\text{ phút}$.\\
\textbf{Đáp án của câu hỏi là D}
}
\end{ex}

\begin{ex} %[5H1H8]
Đổi $2\text{ giờ } 30\text{ phút}$ ra phút ta được
\choice
{$120\text{ phút}$}
{$130\text{ phút}$}
{\True $150\text{ phút}$}
{$230\text{ phút}$}
\loigiai{
Ta có: $2\text{ giờ } 30\text{ phút} = 2 \times 60 + 30 = 120 + 30 = 150\text{ phút}$.\\
\textbf{Đáp án của câu hỏi là C}
}
\end{ex}

\begin{ex} %[5H1H8]
Đổi $4\text{ phút } 5\text{ giây}$ ra giây ta được
\choice
{$45\text{ giây}$}
{\True $245\text{ giây}$}
{$205\text{ giây}$}
{$405\text{ giây}$}
\loigiai{
Ta có: $4\text{ phút } 5\text{ giây} = 4 \times 60 + 5 = 240 + 5 = 245\text{ giây}$.\\
\textbf{Đáp án của câu hỏi là B}
}
\end{ex}

\begin{ex} %[5H1H8]
Đổi $\dfrac{7}{100}\text{ thế kỉ}$ ra năm ta được kết quả là
\choice
{\True $7\text{ năm}$}
{$70\text{ năm}$}
{$700\text{ năm}$}
{$14\text{ năm}$}
\loigiai{
Vì $1\text{ thế kỉ} = 100\text{ năm}$ nên $\dfrac{7}{100}\text{ thế kỉ} = 100 \times \dfrac{7}{100} = 7\text{ năm}$.\\
\textbf{Đáp án của câu hỏi là A}
}
\end{ex}

\begin{ex} %[5H1V8]
Bác Năm thu hoạch được $1\text{ tấn } 250\text{ kg}$ cam. Biết rằng số cam loại I chiếm $\dfrac{3}{10}$ tổng số cam thu hoạch. Số ki-lô-gam cam loại I là
\choice
{$250\text{ kg}$}
{$450\text{ kg}$}
{\True $375\text{ kg}$}
{$875\text{ kg}$}
\loigiai{
Đổi $1\text{ tấn } 250\text{ kg} = 1\,250\text{ kg}$.\\
Số cam loại I là: $1\,250 \times \dfrac{3}{10} = 375\text{ (kg)}$.\\
\textbf{Đáp án của câu hỏi là C}
}
\end{ex}

\begin{ex} %[5H1V8]
Dì Sáu có một mảnh đất hình chữ nhật có chiều dài $12\text{ m}$, chiều rộng $7\text{ m}$. Diện tích của mảnh đất đó là
\choice
{\True $84\text{ m}^2$}
{$38\text{ m}^2$}
{$72\text{ m}^2$}
{$96\text{ m}^2$}
\loigiai{
Diện tích mảnh đất hình chữ nhật là: $12 \times 7 = 84\text{ (m}^2\text{)}$.\\
\textbf{Đáp án của câu hỏi là A}
}
\end{ex}

\begin{ex} %[5H1V8]
Dì Sáu chia đều mảnh đất hình chữ nhật diện tích $84\text{ m}^2$ thành $7$ ô đất bằng nhau để xây phòng trọ. Diện tích mỗi phòng trọ là
\choice
{$7\text{ m}^2$}
{\True $12\text{ m}^2$}
{$10\text{ m}^2$}
{$14\text{ m}^2$}
\loigiai{
Diện tích mỗi ô đất là: $84 : 7 = 12\text{ (m}^2\text{)}$.\\
\textbf{Đáp án của câu hỏi là B}
}
\end{ex}

\textbf{Phần 2. Trắc nghiệm đúng sai}
\setcounterpageref{ex}{0}

\begin{ex} %[5H1H8]
Xét tính đúng sai của các khẳng định sau về đổi đơn vị đo:
\choiceTF
{\True $5\text{ tấn} = 50\text{ tạ}$.}
{$9\text{ tạ} = 900\text{ yến}$.}
{\True $7\text{ thế kỉ} = 700\text{ năm}$.}
{$\dfrac{1}{10}\text{ giờ} = 10\text{ phút}$.}
\loigiai{
\begin{itemize}
    \item $5\text{ tấn} = 5 \times 10 = 50\text{ tạ}$.
    \item $9\text{ tạ} = 9 \times 10 = 90\text{ yến} \neq 900\text{ yến}$.
    \item $7\text{ thế kỉ} = 7 \times 100 = 700\text{ năm}$.
    \item $\dfrac{1}{10}\text{ giờ} = 60 : 10 = 6\text{ phút} \neq 10\text{ phút}$.
\end{itemize}
\begin{itemize}
    \item a) \textbf{Đúng}.
    \item b) \textbf{Sai}: $9\text{ tạ} = 90\text{ yến}$.
    \item c) \textbf{Đúng}.
    \item d) \textbf{Sai}: $\dfrac{1}{10}\text{ giờ} = 6\text{ phút}$.
\end{itemize}
}
\end{ex}

\begin{ex} %[5H1V8]
Bác Năm thu hoạch được $1\text{ tấn } 250\text{ kg}$ cam gồm hai loại: cam loại I và cam loại II. Biết cam loại I chiếm $\dfrac{3}{10}$ tổng số cam thu hoạch. Xét tính đúng sai của các khẳng định sau:
\choiceTF
{\True Tổng khối lượng cam bác Năm thu hoạch được là $1\,250\text{ kg}$.}
{Khối lượng cam loại I là $250\text{ kg}$.}
{\True Khối lượng cam loại II là $875\text{ kg}$.}
{\True Khối lượng cam loại II nhiều hơn khối lượng cam loại I là $500\text{ kg}$.}
\loigiai{
\begin{itemize}
    \item Đổi $1\text{ tấn } 250\text{ kg} = 1\,250\text{ kg}$.
    \item Số cam loại I là: $1\,250 \times \dfrac{3}{10} = 375\text{ kg} \neq 250\text{ kg}$.
    \item Số cam loại II là: $1\,250 - 375 = 875\text{ kg}$.
    \item Hiệu giữa cam loại II và loại I là: $875 - 375 = 500\text{ kg}$.
\end{itemize}
\begin{itemize}
    \item a) \textbf{Đúng}.
    \item b) \textbf{Sai}: Khối lượng cam loại I là $375\text{ kg}$.
    \item c) \textbf{Đúng}.
    \item d) \textbf{Đúng}.
\end{itemize}
}
\end{ex}

\textbf{Phần 3. Trả lời ngắn}
\setcounterpageref{ex}{0}

\begin{ex} %[5H1H8]
Đổi $9\text{ tạ}$ ra yến. Kết quả là bao nhiêu yến?
\shortans{90}
\loigiai{
\textbf{Đáp số:} $90$\\
$9\text{ tạ} = 9 \times 10 = 90\text{ yến}$.
}
\end{ex}

\begin{ex} %[5H1H8]
Đổi $4\text{ phút } 5\text{ giây}$ ra đơn vị giây.
\shortans{245}
\loigiai{
\textbf{Đáp số:} $245$\\
$4\text{ phút } 5\text{ giây} = 4 \times 60 + 5 = 245\text{ giây}$.
}
\end{ex}

\begin{ex} %[5H1H8]
Trong các góc có số đo lần lượt là $60^\circ, 90^\circ, 120^\circ, 45^\circ, 150^\circ$, có bao nhiêu góc là góc tù?
\shortans{2}
\loigiai{
\textbf{Đáp số:} $2$\\
Các góc tù (lớn hơn $90^\circ$ và nhỏ hơn $180^\circ$) là: $120^\circ$ và $150^\circ$. Có tất cả $2$ góc tù.
}
\end{ex}

\begin{ex} %[5H1V8]
Bác Năm thu hoạch được $1\text{ tấn } 250\text{ kg}$ cam. Biết số cam loại I chiếm $\dfrac{3}{10}$ tổng số cam, còn lại là cam loại II. Hỏi bác Năm thu hoạch được bao nhiêu ki-lô-gam cam loại II?
\shortans{875}
\loigiai{
\textbf{Đáp số:} $875$\\
Tổng số cam là $1\,250\text{ kg}$. Số cam loại I là $1\,250 \times \dfrac{3}{10} = 375\text{ kg}$.\\
Số cam loại II là: $1\,250 - 375 = 875\text{ kg}$.
}
\end{ex}

\begin{ex} %[5H1V8]
Mảnh đất hình chữ nhật có chiều dài $12\text{ m}$ và chiều rộng $7\text{ m}$ được chia thành $7$ ô đất nhỏ có diện tích bằng nhau dạng hình chữ nhật có chiều dài $4\text{ m}$. Hỏi chiều rộng của mỗi ô đất nhỏ là bao nhiêu mét?
\shortans{3}
\loigiai{
\textbf{Đáp số:} $3$\\
Diện tích mỗi ô nhỏ là: $(12 \times 7) : 7 = 12\text{ (m}^2\text{)}$.\\
Chiều rộng mỗi ô nhỏ là: $12 : 4 = 3\text{ (m)}$.
}
\end{ex}

\begin{ex} %[5H1H8]
Đổi $\dfrac{7}{100}\text{ thế kỉ}$ ra đơn vị năm.
\shortans{7}
\loigiai{
\textbf{Đáp số:} $7$\\
$1\text{ thế kỉ} = 100\text{ năm} \Rightarrow \dfrac{7}{100} \times 100 = 7\text{ năm}$.
}
\end{ex}
